
\subsection*{4.7.}

On
revient à la situation de 3.7, où on avait le diagramme commutatif :
@@M0@@
où @@M1@@ est une immersion régulière, @@M2@@ le @@M3@@-schéma obtenu en faisant éclater @@M4@@, et @@M5@@. On suppose @@M6@@ noethérien avec un faisceau ample, et @@M7@@ régulier. On a vu en 3.7 qu'on a alors la suite exacte :
@@M8@@
Nous allons montrer que cette suite exacte permet de décrire toute la situation formée par les @@M9@@ de @@M10@@ à partir de la donnée des @@M11@@-anneaux @@M12@@, @@M13@@, des applications @@M14@@ et @@M15@@, et de l'élément @@M16@@ de @@M17@@. (du moins à torsion près, jusqu'à nouvel ordre).

Soient donc @@M18@@ deux @@M19@@-anneaux, @@M20@@ un homomorphisme additif, @@M21@@ un @@M22@@-homomorphisme, @@M23@@ un entier, et @@M24@@ un élément de @@M25@@ tel que @@M26@@ pour @@M27@@ et que @@M28@@ soit inversible. On suppose de plus vraies les propriétés suivantes :
\begin{enumerate}[label=\roman*)]
\item @@M29@@ est @@M30@@-linéaire (« formule de projection »);
\item @@M31@@ est l'homomorphisme « multiplication par @@M32@@ »;
\item pour tout @@M33@@, et tout @@M34@@, @@M35@@ (relation (4.3.1)).
\end{enumerate}

On utilise d'abord VI.1.1 pour construire une @@M36@@-@@M37@@-algèbre @@M38@@ qui joue le rôle de @@M39@@. Soit @@M40@@ une indéterminée, et @@M41@@ la @@M42@@-@@M43@@-algèbre engendrée par le générateur @@M44@@ soumis aux relations @@M45@@ pour @@M46@@ (V 4.9). Soit :
@@M47@@
